\iftestbuild % TEST-ONLY-BEGIN
\setcounter{section}{1}\setcounter{theorem}{25}
\fi % TEST-ONLY-END
% Source: book pp. 39--43 (PDF 53--57); solutions pp. 35--41.
% A numbered display (Axler's 2.29, 2.44, 2.45) takes its number from the
% shared item counter: \axeqnum{2.29} inside equation* (same definition in 2A, 2C).
\providecommand\axeqnum[1]{\refstepcounter{theorem}\tag*{\thetheorem}\label{ax:#1}}
\section{Bases}

In the previous section, we discussed linearly independent lists and we also
discussed spanning lists. Now we bring these concepts together by considering
lists that have both properties.

\begin{definition}[Basis]\label{ax:2.26}
A \emph{basis} of $V$ is a list of vectors in $V$ that is linearly independent
and spans $V$.
\end{definition}

\begin{example}[Bases]\label{ax:2.27}\leavevmode
\begin{parts}
\item The list $(1,0,\dots,0),(0,1,0,\dots,0),\dots,(0,\dots,0,1)$ is a basis of
  $\F^n$, called the \emph{standard basis} of $\F^n$.
\item The list $(1,2),(3,5)$ is a basis of $\F^2$. Note that this list has
  length two, which is the same as the length of the standard basis of $\F^2$.
  In the next section, we will see that this is not a coincidence.
\item The list $(1,2,-4),(7,-5,6)$ is linearly independent in $\F^3$ but is not
  a basis of $\F^3$ because it does not span $\F^3$.
\item The list $(1,2),(3,5),(4,13)$ spans $\F^2$ but is not a basis of $\F^2$
  because it is not linearly independent.
\item The list $(1,1,0),(0,0,1)$ is a basis of $\{(x,x,y)\in\F^3 : x,y\in\F\}$.
\item The list $(1,-1,0),(1,0,-1)$ is a basis of
  \[ \{(x,y,z)\in\F^3 : x+y+z=0\}. \]
\item The list $1,z,\dots,z^m$ is a basis of $\Poly_m(\F)$, called the
  \emph{standard basis} of $\Poly_m(\F)$.
\end{parts}
\end{example}

In addition to the standard basis, $\F^n$ has many other bases. For example,
\[ (7,5),(-4,9) \quad\text{and}\quad (1,2),(3,5) \]
are both bases of $\F^2$.

The next result helps explain why bases are useful. Recall that ``uniquely''
means ``in only one way''.

\begin{result}[Criterion for basis]\label{ax:2.28}
A list $v_1,\dots,v_n$ of vectors in $V$ is a basis of $V$ if and only if every
$v\in V$ can be written uniquely in the form
\begin{equation*}
v = a_1v_1+\cdots+a_nv_n, \axeqnum{2.29}
\end{equation*}
where $a_1,\dots,a_n\in\F$.
\end{result}

\begin{proof}
\marginnote[2em]{This proof is essentially a repetition of the ideas that led us to
the definition of linear independence.}%
First suppose that $v_1,\dots,v_n$ is a basis of $V$. Let $v\in V$. Because
$v_1,\dots,v_n$ spans $V$, there exist $a_1,\dots,a_n\in\F$ such that
\ref{ax:2.29} holds. To show that the representation in \ref{ax:2.29} is
unique, suppose $c_1,\dots,c_n$ are scalars such that we also have
\[ v = c_1v_1+\cdots+c_nv_n. \]
Subtracting the last equation from \ref{ax:2.29}, we get
\[ 0 = (a_1-c_1)v_1+\cdots+(a_n-c_n)v_n. \]
This implies that each $a_k-c_k$ equals $0$ (because $v_1,\dots,v_n$ is linearly
independent). Hence $a_1=c_1,\dots,a_n=c_n$. We have the desired uniqueness,
completing the proof in one direction.

For the other direction, suppose every $v\in V$ can be written uniquely in the
form given by \ref{ax:2.29}. This implies that the list $v_1,\dots,v_n$ spans
$V$. To show that $v_1,\dots,v_n$ is linearly independent, suppose
$a_1,\dots,a_n\in\F$ are such that
\[ 0 = a_1v_1+\cdots+a_nv_n. \]
The uniqueness of the representation \ref{ax:2.29} (taking $v=0$) now implies
that $a_1=\cdots=a_n=0$. Thus $v_1,\dots,v_n$ is linearly independent and hence
is a basis of $V$.
\end{proof}

A spanning list in a vector space may not be a basis because it is not linearly
independent. Our next result says that given any spanning list, some (possibly
none) of the vectors in it can be discarded so that the remaining list is
linearly independent and still spans the vector space.

As an example in the vector space $\F^2$, if the procedure in the proof below is
applied to the list $(1,2),(3,6),(4,7),(5,9)$, then the second and fourth
vectors will be removed. This leaves $(1,2),(4,7)$, which is a basis of $\F^2$.

\begin{result}[Every spanning list contains a basis]\label{ax:2.30}
Every spanning list in a vector space can be reduced to a basis of the vector
space.
\end{result}

\begin{proof}
Suppose $v_1,\dots,v_n$ spans $V$. We want to remove some of the vectors from
$v_1,\dots,v_n$ so that the remaining vectors form a basis of $V$. We do this
through the multistep process described below.

Start with $B$ equal to the list $v_1,\dots,v_n$.
\begin{description}[style=nextline,leftmargin=1em,itemsep=2pt,parsep=3pt]
\item[Step 1] If $v_1=0$, then delete $v_1$ from $B$. If $v_1\ne0$, then leave
  $B$ unchanged.
\item[Step $\bm{k}$] If $v_k$ is in $\spn(v_1,\dots,v_{k-1})$, then
  delete $v_k$ from the list $B$. If $v_k$ is not in $\spn(v_1,\dots,v_{k-1})$,
  then leave $B$ unchanged.
\end{description}
Stop the process after step $n$, getting a list $B$. This list $B$ spans $V$
because our original list spanned $V$ and we have discarded only vectors that
were already in the span of the previous vectors. The process ensures that no
vector in $B$ is in the span of the previous ones. Thus $B$ is linearly
independent, by the linear dependence lemma (\ref{ax:2.19}). Hence $B$ is a
basis of $V$.
\end{proof}

We now come to an important corollary of the previous result.

\begin{result}[Basis of finite-dimensional vector space]\label{ax:2.31}
Every finite-dimensional vector space has a basis.
\end{result}

\begin{proof}
By definition, a finite-dimensional vector space has a spanning list. The
previous result tells us that each spanning list can be reduced to a basis.
\end{proof}

Our next result is in some sense a dual of \ref{ax:2.30}, which said that every
spanning list can be reduced to a basis. Now we show that given any linearly
independent list, we can adjoin some additional vectors (this includes the
possibility of adjoining no additional vectors) so that the extended list is
still linearly independent but also spans the space.

\begin{result}[Every linearly independent list extends to a basis]\label{ax:2.32}
Every linearly independent list of vectors in a finite-dimensional vector space
can be extended to a basis of the vector space.
\end{result}

\begin{proof}
Suppose $u_1,\dots,u_m$ is linearly independent in a finite-dimensional vector
space $V$. Let $w_1,\dots,w_n$ be a list of vectors in $V$ that spans $V$. Thus
the list
\[ u_1,\dots,u_m,w_1,\dots,w_n \]
spans $V$. Applying the procedure of the proof of \ref{ax:2.30} to reduce this
list to a basis of $V$ produces a basis consisting of the vectors
$u_1,\dots,u_m$ and some of the $w$'s (none of the $u$'s get deleted in this
procedure because $u_1,\dots,u_m$ is linearly independent).
\end{proof}

As an example in $\F^3$, suppose we start with the linearly independent list
$(2,3,4),(9,6,8)$. If we take $w_1,w_2,w_3$ to be the standard basis of $\F^3$,
then applying the procedure in the proof above produces the list
\[ (2,3,4),(9,6,8),(0,1,0), \]
which is a basis of $\F^3$.

\marginnote{Using the same ideas but more advanced tools, the next result can be
proved without the hypothesis that $V$ is finite-dimensional.}%
As an application of the result above, we now show that every subspace of a
finite-dimensional vector space can be paired with another subspace to form a
direct sum of the whole space.

\begin{result}[Every subspace of $V$ is part of a direct sum equal to
$V$]\label{ax:2.33}
Suppose $V$ is finite-dimensional and $U$ is a subspace of $V$. Then there is a
subspace $W$ of $V$ such that $V=U\oplus W$.
\end{result}

\begin{proof}
Because $V$ is finite-dimensional, so is $U$ (see \ref{ax:2.25}). Thus there is
a basis $u_1,\dots,u_m$ of $U$ (by \ref{ax:2.31}). Of course $u_1,\dots,u_m$ is a
linearly independent list of vectors in $V$. Hence this list can be extended to
a basis $u_1,\dots,u_m,w_1,\dots,w_n$ of $V$ (by \ref{ax:2.32}). Let
$W=\spn(w_1,\dots,w_n)$.

To prove that $V=U\oplus W$, by \ref{ax:1.46} we only need to show that
\[ V=U+W \quad\text{and}\quad U\cap W=\{0\}. \]

To prove the first equation above, suppose $v\in V$. Then, because the list
$u_1,\dots,u_m,w_1,\dots,w_n$ spans $V$, there exist
$a_1,\dots,a_m,b_1,\dots,b_n\in\F$ such that
\[ v = \underbrace{a_1u_1+\cdots+a_mu_m}_{u}
     + \underbrace{b_1w_1+\cdots+b_nw_n}_{w}. \]
We have $v=u+w$, where $u\in U$ and $w\in W$ are defined as above. Thus
$v\in U+W$, completing the proof that $V=U+W$.

To show that $U\cap W=\{0\}$, suppose $v\in U\cap W$. Then there exist scalars
$a_1,\dots,a_m,b_1,\dots,b_n\in\F$ such that
\[ v = a_1u_1+\cdots+a_mu_m = b_1w_1+\cdots+b_nw_n. \]
Thus
\[ a_1u_1+\cdots+a_mu_m-b_1w_1-\cdots-b_nw_n = 0. \]
Because $u_1,\dots,u_m,w_1,\dots,w_n$ is linearly independent, this implies that
\[ a_1=\cdots=a_m=b_1=\cdots=b_n=0. \]
Thus $v=0$, completing the proof that $U\cap W=\{0\}$.
\end{proof}

% ------------------------------------------------------------------ exercises
\exercisesheading{2B}

\begin{exercise}
Find all vector spaces that have exactly one basis.
\end{exercise}
\begin{solution}
\solnote{Manual Remark}{The proof of this exercise fails if $\F$ contains only
two elements. In this case, an additional vector space with exactly one basis is
the vector space of dimension one over $\F$, since the only nonzero scalar is
$1$, and any nonzero vector $v$ in the vector space is a basis.}

Since an infinite-dimensional vector space has infinitely many bases, we only
need to consider finite-dimensional vector spaces. Let $V$ be a
finite-dimensional vector space, and let $v_1,\dots,v_n$ be a basis of $V$. Then
any other basis of $V$ must have the same length as $v_1,\dots,v_n$, which is
$n$. If $n\ge1$, consider $\alpha\in\F$ such that $\alpha\ne0$ and $\alpha\ne1$.
Then the list $\alpha v_1,v_2,\dots,v_n$ still spans $V$, since we may recover
$v_1=\alpha^{-1}(\alpha v_1)$. Moreover, the list $\alpha v_1,v_2,\dots,v_n$ is
linearly independent, since if
\[ \beta_1(\alpha v_1)+\beta_2v_2+\cdots+\beta_nv_n = 0, \]
then we may write
\[ (\alpha\beta_1)v_1+\beta_2v_2+\cdots+\beta_nv_n = 0. \]
Since $v_1,\dots,v_n$ is linearly independent, we must have $\alpha\beta_1=0$,
which implies that $\beta_1=0$. Then we also have $\beta_2=\cdots=\beta_n=0$.
Thus, the list $\alpha v_1,v_2,\dots,v_n$ is linearly independent. Therefore, if
$n\ge1$, then $V$ has more than one basis. If $n=0$, then $V=\{0\}$. The empty
list spans $V$, and the empty list must be linearly independent, since there
are no nonzero scalars to consider. Thus, the empty list is a basis of $V$.
Moreover, this is the only basis of $V$, since $0$ is linearly dependent.
Therefore, the only vector space that has exactly one basis is the zero vector
space.
\end{solution}

\begin{exercise}
Verify all assertions in Example~\ref{ax:2.27}.
\end{exercise}
\begin{solution}\leavevmode
\begin{parts}
\item Suppose $\alpha_1,\dots,\alpha_n\in\F$ such that
  \[ \alpha_1(1,0,\dots,0)+\alpha_2(0,1,0,\dots,0)+\cdots+\alpha_n(0,\dots,0,1)
     = 0. \]
  Then each $\alpha_i=0$, so the list is linearly independent. Moreover, any
  vector $(x_1,\dots,x_n)\in\F^n$ can be written as
  \[ (x_1,\dots,x_n) = x_1(1,0,\dots,0)+x_2(0,1,0,\dots,0)+\cdots
     +x_n(0,\dots,0,1), \]
  so the list spans $\F^n$. Thus, the list is a basis of $\F^n$.
\item If $(3,5)=a(1,2)$ for some $a\in\F$, then we must have $3=a$ and $5=2a$,
  which is impossible. Thus, the list is linearly independent. Moreover, for
  any element $(x,y)\in\F^2$, we can solve the system of equations
  \begin{align*}
  x &= 3\alpha+\beta,\\
  y &= 5\alpha+2\beta
  \end{align*}
  for $\alpha,\beta\in\F$. Solving this system of equations, we find that
  \[ \alpha=2x-y, \quad \beta=-5x+3y. \]
  Since we can find $\alpha,\beta\in\F$ for any $(x,y)\in\F^2$, the list spans
  $\F^2$. Thus, the list is a basis of $\F^2$.
\item Note that $\F^3$ is spanned by the linearly independent list
  $(1,0,0),\allowbreak(0,1,0),\allowbreak(0,0,1)$. Since the length of a spanning list must be at
  least the length of a linearly independent list, the list
  $(1,2,-4),(7,-5,6)$ cannot span $\F^3$.
\item From part (b), we saw that $(1,2),(3,5)$ is a basis of $\F^2$. Since
  $(4,13)=-5(3,5)+19(1,2)$, the list $(1,2),(3,5),(4,13)$ is linearly
  dependent. Moreover, since $(1,2),(3,5)$ is a basis of $\F^2$, the list
  $(1,2),(3,5),(4,13)$ spans $\F^2$. Thus, the list is not a basis of $\F^2$.
\item Let $U$ be the subspace in question. If $\alpha,\beta\in\F$ such that
  \[ \alpha(1,1,0)+\beta(0,0,1) = (0,0,0), \]
  then we must have $\alpha=\beta=0$, so the list is linearly independent.
  Moreover, for any $(x,x,y)\in U$, we can write
  \[ (x,x,y) = x(1,1,0)+y(0,0,1), \]
  so the list spans $U$. Thus, the list is a basis of $U$.
\item Let $U$ be the subspace in question. If $\alpha,\beta\in\F$ such that
  \[ \alpha(1,-1,0)+\beta(1,0,-1) = (0,0,0), \]
  then we must have $\alpha=\beta=0$, so the list is linearly independent.
  Moreover, for any $(x,y,z)\in U$, note that $x+y+z=0$, implying that
  $x=-y-z$. Then we can write
  \[ (x,y,z) = (-y-z,y,z) = -y(1,-1,0)-z(1,0,-1), \]
  so the list spans $U$. Thus, the list is a basis of $U$.
\item If $\alpha_0,\dots,\alpha_m\in\F$ such that
  \[ \alpha_0+\alpha_1z+\cdots+\alpha_mz^m = 0, \]
  then we must have $\alpha_0=\cdots=\alpha_m=0$, so the list is linearly
  independent. Moreover, any polynomial $p(z)\in\Poly_m(\F)$ can be written as
  \[ p(z) = \alpha_0+\alpha_1z+\cdots+\alpha_mz^m, \]
  so the list spans $\Poly_m(\F)$. Thus, the list is a basis of $\Poly_m(\F)$.
\end{parts}
\end{solution}

\begin{exercise}\leavevmode
\begin{parts}
\item Let $U$ be the subspace of $\R^5$ defined by
  \[ U = \{(x_1,x_2,x_3,x_4,x_5)\in\R^5 : x_1=3x_2 \text{ and } x_3=7x_4\}. \]
  Find a basis of $U$.
\item Extend the basis in (a) to a basis of $\R^5$.
\item Find a subspace $W$ of $\R^5$ such that $\R^5=U\oplus W$.
\end{parts}
\end{exercise}
\begin{solution}\leavevmode
\begin{parts}
\item Note that we can rewrite the subspace $U$ as
  \[ U = \{(3x_2,x_2,7x_4,x_4,x_5)\in\R^5 \mid x_2,x_4,x_5\in\R\}. \]
  Then we can write any vector in $U$ as a linear combination of the following
  three vectors:
  \[ u_1=(3,1,0,0,0), \quad u_2=(0,0,7,1,0), \quad u_3=(0,0,0,0,1). \]
  Moreover, these three vectors are linearly independent, since the only way to
  write the zero vector as a linear combination of $u_1,u_2,u_3$ is to take all
  coefficients to be zero. Thus, $\{u_1,u_2,u_3\}$ is a basis of $U$.
\item We can extend the basis $\{u_1,u_2,u_3\}$ of $U$ to a basis of $\R^5$ by
  adding two more vectors that are linearly independent of $u_1,u_2,u_3$. For
  example, we can add the vectors
  \[ u_4=(1,0,0,0,0), \quad u_5=(0,0,1,0,0), \]
  which are linearly independent of $u_1,u_2,u_3$: since $u_1$ is the only
  vector in the list with a nonzero first component, $u_4$ cannot be a linear
  combination of $u_1,u_2,u_3$. Similarly, since $u_2$ is the only vector in the
  list with a nonzero third component, $u_5$ cannot be a linear combination of
  $u_1,u_2,u_3,u_4$. Thus, $\{u_1,u_2,u_3,u_4,u_5\}$ is a basis of $\R^5$.
\item Let $W$ be the subspace of $\R^5$ spanned by the vectors $u_4$ and $u_5$.
  Then we have
  \[ \R^5 = U\oplus W, \]
  since $U\cap W=\{0\}$ and every vector in $\R^5$ can be written as a sum of a
  vector in $U$ and a vector in $W$.
\end{parts}
\end{solution}

\begin{exercise}\leavevmode
\begin{parts}
\item Let $U$ be the subspace of $\C^5$ defined by
  \[ U = \{(z_1,z_2,z_3,z_4,z_5)\in\C^5 : 6z_1=z_2 \text{ and }
     z_3+2z_4+3z_5=0\}. \]
  Find a basis of $U$.
\item Extend the basis in (a) to a basis of $\C^5$.
\item Find a subspace $W$ of $\C^5$ such that $\C^5=U\oplus W$.
\end{parts}
\end{exercise}
\begin{solution}\leavevmode
\begin{parts}
\item Note that we can rewrite the subspace $U$ as
  \[ U = \{(z_1,6z_1,-2z_4-3z_5,z_4,z_5)\in\C^5 \mid z_1,z_4,z_5\in\C\}. \]
  Then we can write any vector in $U$ as a linear combination of the following
  three vectors:
  \[ u_1=(1,6,0,0,0), \quad u_2=(0,0,-2,1,0), \quad u_3=(0,0,-3,0,1). \]
  Moreover, these three vectors are linearly independent, since the only way to
  write the zero vector as a linear combination of $u_1,u_2,u_3$ is to take all
  coefficients to be zero. Thus, $\{u_1,u_2,u_3\}$ is a basis of $U$.
\item We can extend the basis $\{u_1,u_2,u_3\}$ of $U$ to a basis of $\C^5$ by
  adding two more vectors that are linearly independent of $u_1,u_2,u_3$. For
  example, we can add the vectors
  \[ u_4=(1,0,0,0,0), \quad u_5=(0,0,1,0,0), \]
  which are linearly independent of $u_1,u_2,u_3$: since $u_1$ is the only
  vector in the list with a nonzero first component, $u_4$ cannot be a linear
  combination of $u_1,u_2,u_3$. Similarly, $u_2$ and $u_3$ have nonzero third
  components, but any linear combination of $u_2$ and $u_3$ will have nonzero
  fourth and fifth components, so $u_5$ cannot be a linear combination of
  $u_1,u_2,u_3,u_4$. Thus, $\{u_1,u_2,u_3,u_4,u_5\}$ is a basis of $\C^5$.
\item Let $W$ be the subspace of $\C^5$ spanned by the vectors $u_4$ and $u_5$.
  Then we have
  \[ \C^5 = U\oplus W, \]
  since $U\cap W=\{0\}$ and every vector in $\C^5$ can be written as a sum of a
  vector in $U$ and a vector in $W$.
\end{parts}
\end{solution}

\begin{exercise}
Suppose $V$ is finite-dimensional and $U,W$ are subspaces of $V$ such that
$V=U+W$. Prove that there exists a basis of $V$ consisting of vectors in
$U\cup W$.
\end{exercise}
\begin{solution}
Let $u_1,\dots,u_m$ be a basis of $U$ and $w_1,\dots,w_n$ a basis of $W$. Then
the list $u_1,\dots,u_m,w_1,\dots,w_n$ spans $V$, since any vector in $V$ can be
written as a sum of a vector in $U$ and a vector in $W$. We can extract a basis
of $V$ from this list by removing any linearly dependent vectors. Thus, there
exists a basis of $V$ consisting of vectors in $U\cup W$.
\end{solution}

\begin{exercise}
Prove or give a counterexample: If $p_0,p_1,p_2,p_3$ is a list in $\Poly_3(\F)$
such that none of the polynomials $p_0,p_1,p_2,p_3$ has degree $2$, then
$p_0,p_1,p_2,p_3$ is not a basis of $\Poly_3(\F)$.
\end{exercise}
\begin{solution}
This is not true. Let $p_0=1$, $p_1=x$, $p_2=x^3$, $p_3=x^2+x^3$. Then none of
the polynomials $p_0,p_1,p_2,p_3$ has degree $2$. Moreover, the list
$p_0,p_1,p_2,p_3$ is linearly independent, since the only way to write the zero
polynomial as a linear combination of $p_0,p_1,p_2,p_3$ is to take all
coefficients to be zero. Finally, noting that $x^2=p_3-p_2$, then for any
polynomial $p(x)=a_0+a_1x+a_2x^2+a_3x^3\in\Poly_3(\F)$, we can write
\[ p(x) = a_0p_0+a_1p_1+a_2(p_3-p_2)+a_3p_2
        = a_0p_0+a_1p_1+(a_3-a_2)p_2+a_2p_3, \]
which shows that the list spans $\Poly_3(\F)$. Thus, the list
$p_0,p_1,p_2,p_3$ is a basis of $\Poly_3(\F)$.
\end{solution}

\begin{exercise}
Suppose $v_1,v_2,v_3,v_4$ is a basis of $V$. Prove that
\[ v_1+v_2,\,v_2+v_3,\,v_3+v_4,\,v_4 \]
is also a basis of $V$.
\end{exercise}
\begin{solution}
\solnote{Manual Remark}{The technique of transforming the scalars
$\beta_1,\beta_2,\beta_3,\beta_4$ into the scalars
$\alpha_1,\alpha_2,\alpha_3,\alpha_4$ is good to remember. You will later learn
that this is called a \textup{change of basis}.}

Let $\alpha_1,\dots,\alpha_4\in\F$ such that
\[ \alpha_1(v_1+v_2)+\alpha_2(v_2+v_3)+\alpha_3(v_3+v_4)+\alpha_4v_4 = 0. \]
Then we may write
\[ \alpha_1v_1+(\alpha_1+\alpha_2)v_2+(\alpha_2+\alpha_3)v_3
   +(\alpha_3+\alpha_4)v_4 = 0. \]
Since $v_1,v_2,v_3,v_4$ is a basis of $V$, we must have
\[ \alpha_1=0, \quad \alpha_1+\alpha_2=0, \quad \alpha_2+\alpha_3=0, \quad
   \alpha_3+\alpha_4=0. \]
From the first equation, we have $\alpha_1=0$. Substituting this into the second
equation gives us $\alpha_2=0$. Substituting this into the third equation gives
us $\alpha_3=0$. Finally, substituting this into the fourth equation gives us
$\alpha_4=0$. Thus, the list is linearly independent. Moreover, since
$v_1,v_2,v_3,v_4$ is a basis of $V$, then any vector $v\in V$ can be written as
\[ v = \beta_1v_1+\beta_2v_2+\beta_3v_3+\beta_4v_4 \]
for some $\beta_1,\dots,\beta_4\in\F$. Using the expressions from proving the
linear independence of the list, we may show that there exist
$\alpha_1,\dots,\alpha_4\in\F$ such that
\[ \begin{cases}
   \beta_1=\alpha_1,\\
   \beta_2=\alpha_1+\alpha_2,\\
   \beta_3=\alpha_2+\alpha_3,\\
   \beta_4=\alpha_3+\alpha_4.
   \end{cases} \]
Then $v\in V$ may be written as
\[ v = \alpha_1(v_1+v_2)+\alpha_2(v_2+v_3)+\alpha_3(v_3+v_4)+\alpha_4v_4, \]
where
\[ \alpha_1=\beta_1, \quad \alpha_2=\beta_2-\beta_1, \quad
   \alpha_3=\beta_3-\beta_2+\beta_1, \quad
   \alpha_4=\beta_4-\beta_3+\beta_2-\beta_1. \]
Thus, the list spans $V$. Therefore, the list is a basis of $V$.
\end{solution}

\begin{exercise}
Prove or give a counterexample: If $v_1,v_2,v_3,v_4$ is a basis of $V$ and $U$ is
a subspace of $V$ such that $v_1,v_2\in U$ and $v_3\notin U$ and $v_4\notin U$,
then $v_1,v_2$ is a basis of $U$.
\end{exercise}
\begin{solution}
This is not true. Let $V=\R^4$ and $v_1,v_2,v_3,v_4$ be the standard basis of
$\R^4$. Let $U$ be the subspace of $\R^4$ defined by
\[ U = \spn(v_1,v_2,v_3+v_4) = \{(x,y,z,z)\in\R^4 \mid x,y,z\in\R\}. \]
Then $v_1,v_2\in U$ and $v_3,v_4\notin U$. However, $v_1,v_2$ does not span
$U$, since $(0,0,1,1)\in U$ but cannot be written as a linear combination of
$v_1$ and $v_2$. Thus, $v_1,v_2$ is not a basis of $U$.
\end{solution}

\begin{exercise}
Suppose $v_1,\dots,v_m$ is a list of vectors in $V$. For $k\in\{1,\dots,m\}$, let
\[ w_k = v_1+\cdots+v_k. \]
Show that $v_1,\dots,v_m$ is a basis of $V$ if and only if $w_1,\dots,w_m$ is a
basis of $V$.
\end{exercise}
\begin{solution}
By Exercise~2A14, we know that $v_1,\dots,v_m$ is linearly independent if and
only if $w_1,\dots,w_m$ is linearly independent. By Exercise~2A3, we know that
$v_1,\dots,v_m$ spans $V$ if and only if $w_1,\dots,w_m$ spans $V$. Thus,
$v_1,\dots,v_m$ is a basis of $V$ if and only if $w_1,\dots,w_m$ is a basis
of $V$.
\end{solution}

\begin{exercise}
Suppose $U$ and $W$ are subspaces of $V$ such that $V=U\oplus W$. Suppose also
that $u_1,\dots,u_m$ is a basis of $U$ and $w_1,\dots,w_n$ is a basis of $W$.
Prove that
\[ u_1,\dots,u_m,w_1,\dots,w_n \]
is a basis of $V$.
\end{exercise}
\begin{solution}
Let $u_1,\dots,u_m$ be a basis of $U$ and $w_1,\dots,w_n$ be a basis of $W$. We
claim that the list
\[ u_1,\dots,u_m,w_1,\dots,w_n \]
is a basis of $V$. Let $\alpha_1,\dots,\alpha_m,\beta_1,\dots,\beta_n\in\F$ such
that
\[ \alpha_1u_1+\cdots+\alpha_mu_m+\beta_1w_1+\cdots+\beta_nw_n = 0. \]
Then we may write
\[ \alpha_1u_1+\cdots+\alpha_mu_m = -\beta_1w_1-\cdots-\beta_nw_n. \]
Since the left-hand side is in $U$ and the right-hand side is in $W$, and
$U\cap W=\{0\}$, we must have
\[ \alpha_1u_1+\cdots+\alpha_mu_m = 0 \quad\text{and}\quad
   \beta_1w_1+\cdots+\beta_nw_n = 0. \]
Since $u_1,\dots,u_m$ is a basis of $U$, we have $\alpha_1=\cdots=\alpha_m=0$.
Similarly, since $w_1,\dots,w_n$ is a basis of $W$, we have
$\beta_1=\cdots=\beta_n=0$. Thus, the list is linearly independent. Moreover,
since $V=U\oplus W$, then $v=u+w$ for some $u\in U$ and $w\in W$. Then we may
write
\[ u = \gamma_1u_1+\cdots+\gamma_mu_m \quad\text{and}\quad
   w = \delta_1w_1+\cdots+\delta_nw_n \]
for some $\gamma_1,\dots,\gamma_m,\delta_1,\dots,\delta_n\in\F$. Hence,
\[ v = u+w = \gamma_1u_1+\cdots+\gamma_mu_m+\delta_1w_1+\cdots+\delta_nw_n, \]
which shows that the list spans $V$. Therefore, the list is a basis of $V$.
\end{solution}

\begin{exercise}
Suppose $V$ is a real vector space. Show that if $v_1,\dots,v_n$ is a basis of
$V$ (as a real vector space), then $v_1,\dots,v_n$ is also a basis of the
complexification $V_{\C}$ (as a complex vector space).
\exnote{See Exercise~8 in Section~1B for the definition of the
complexification $V_{\C}$.}
\end{exercise}
\begin{solution}
Let $v_1,\dots,v_n$ be a basis of $V$. We claim that $v_1,\dots,v_n$ is also a
basis of $V_{\C}$. Let $\alpha_1,\dots,\alpha_n\in\C$ such that
\[ \alpha_1v_1+\cdots+\alpha_nv_n = 0. \]
Then we may write each $\alpha_k=a_k+b_ki$ for some $a_k,b_k\in\R$. Then we have
\[ (a_1+b_1i)v_1+\cdots+(a_n+b_ni)v_n = 0. \]
This implies that
\[ (a_1v_1+\cdots+a_nv_n)+i(b_1v_1+\cdots+b_nv_n) = 0. \]
Since the left-hand side is a sum of two vectors in $V$, and $V$ is a real
vector space, we must have
\[ a_1v_1+\cdots+a_nv_n = 0 \quad\text{and}\quad b_1v_1+\cdots+b_nv_n = 0. \]
Since $v_1,\dots,v_n$ is a basis of $V$, we have $a_1=\cdots=a_n=0$ and
$b_1=\cdots=b_n=0$. Thus, $\alpha_1=\cdots=\alpha_n=0$, which shows that the
list is linearly independent. Moreover, let $v\in V_{\C}$. Then we may write
$v=u+wi$ for some $u,w\in V$. Since $v_1,\dots,v_n$ is a basis of $V$, we may
write
\[ u = \gamma_1v_1+\cdots+\gamma_nv_n \quad\text{and}\quad
   w = \delta_1v_1+\cdots+\delta_nv_n \]
for some $\gamma_1,\dots,\gamma_n,\delta_1,\dots,\delta_n\in\R$. Hence,
\[ v = u+wi = (\gamma_1+\delta_1i)v_1+\cdots+(\gamma_n+\delta_ni)v_n, \]
which shows that the list spans $V_{\C}$. Therefore, the list is a basis of
$V_{\C}$.
\end{solution}
